%==== KNITR ===================================================================%



%==== START ===================================================================%

\documentclass[10pt]{beamer}\usepackage[]{graphicx}\usepackage[table]{xcolor}
% maxwidth is the original width if it is less than linewidth
% otherwise use linewidth (to make sure the graphics do not exceed the margin)
\makeatletter
\def\maxwidth{ %
  \ifdim\Gin@nat@width>\linewidth
    \linewidth
  \else
    \Gin@nat@width
  \fi
}
\makeatother

\definecolor{fgcolor}{rgb}{0.345, 0.345, 0.345}
\newcommand{\hlnum}[1]{\textcolor[rgb]{0.686,0.059,0.569}{#1}}%
\newcommand{\hlsng}[1]{\textcolor[rgb]{0.192,0.494,0.8}{#1}}%
\newcommand{\hlcom}[1]{\textcolor[rgb]{0.678,0.584,0.686}{\textit{#1}}}%
\newcommand{\hlopt}[1]{\textcolor[rgb]{0,0,0}{#1}}%
\newcommand{\hldef}[1]{\textcolor[rgb]{0.345,0.345,0.345}{#1}}%
\newcommand{\hlkwa}[1]{\textcolor[rgb]{0.161,0.373,0.58}{\textbf{#1}}}%
\newcommand{\hlkwb}[1]{\textcolor[rgb]{0.69,0.353,0.396}{#1}}%
\newcommand{\hlkwc}[1]{\textcolor[rgb]{0.333,0.667,0.333}{#1}}%
\newcommand{\hlkwd}[1]{\textcolor[rgb]{0.737,0.353,0.396}{\textbf{#1}}}%
\let\hlipl\hlkwb

\usepackage{framed}
\makeatletter
\newenvironment{kframe}{%
 \def\at@end@of@kframe{}%
 \ifinner\ifhmode%
  \def\at@end@of@kframe{\end{minipage}}%
  \begin{minipage}{\columnwidth}%
 \fi\fi%
 \def\FrameCommand##1{\hskip\@totalleftmargin \hskip-\fboxsep
 \colorbox{shadecolor}{##1}\hskip-\fboxsep
     % There is no \\@totalrightmargin, so:
     \hskip-\linewidth \hskip-\@totalleftmargin \hskip\columnwidth}%
 \MakeFramed {\advance\hsize-\width
   \@totalleftmargin\z@ \linewidth\hsize
   \@setminipage}}%
 {\par\unskip\endMakeFramed%
 \at@end@of@kframe}
\makeatother

\definecolor{shadecolor}{rgb}{.97, .97, .97}
\definecolor{messagecolor}{rgb}{0, 0, 0}
\definecolor{warningcolor}{rgb}{1, 0, 1}
\definecolor{errorcolor}{rgb}{1, 0, 0}
\newenvironment{knitrout}{}{} % an empty environment to be redefined in TeX

\usepackage{alltt}

% --- THEME SETUP ---
\usetheme{Madrid}
\usecolortheme{beaver} % Restores the Red theme

% --- PACKAGES ---
\usepackage[utf8]{inputenc}
\usepackage{amsmath,amssymb}
\usepackage{graphicx}
\usepackage{booktabs}
\usepackage{hyperref}
\usepackage{natbib}
\setcitestyle{authoryear,open={(},close={)}}
\usepackage[table]{xcolor}
\usepackage{caption}

% --- CUSTOM COLOR & FONT OVERRIDES (B) ---
% 1. Define the specific red (Beaver default is usually dark red)
\definecolor{beaverRed}{RGB}{140,45,45} 
\definecolor{MainColour}{RGB}{0,72, 144} 
\definecolor{grey}{RGB}{216,221,227} 
\definecolor{darkgrey}{RGB}{149,153,158} 

% 2. Force Titles to be White on Red
\setbeamercolor{title}{bg=MainColour, fg=white}
\setbeamercolor{frametitle}{bg=MainColour, fg=white}
\setbeamercolor{subtitle}{fg=white} % Subtitle on title page

% 3. Force Footer to be White on Red
\setbeamercolor{date in head/foot}{bg=MainColour, fg=white} 
\setbeamercolor{footline}{bg=MainColour, fg=white}

%%%% Adjusts the text boxes (alterboxes) %%%%
\setbeamercolor{block title alerted}{bg=MainColour, fg=white}
\setbeamercolor{block body alerted}{bg=grey, fg=black}
\setbeamercolor{block title}{bg=darkgrey, fg=white}
\setbeamercolor{block body}{bg=grey, fg=black}

%%%% Enable caption numbering.
\setbeamertemplate{caption}[numbered]

%%%% Customizes the footline / footnoes %%%%
\setbeamertemplate{footline}{%
  \leavevmode%
  \hbox{%
    % LEFT BOX: Date
    \begin{beamercolorbox}[wd=6.4cm,ht=2.25ex,dp=1ex,leftskip=0.5cm]{date in head/foot}%
      \usebeamerfont{date in head/foot}\insertdate
    \end{beamercolorbox}%
    % RIGHT BOX: Page Number
    % 1. REMOVED [right] option to stop it from forcing the text to the edge.
    % 2. ADDED \hfill to push the number to the right side.
    % 3. ADDED \hspace*{...} to create the gap.
    \begin{beamercolorbox}[wd=6.4cm,ht=2.25ex,dp=1ex]{date in head/foot}%
      \usebeamerfont{date in head/foot}\hfill\insertframenumber\hspace*{0.75cm}% <--- Adjust this value
    \end{beamercolorbox}}%
  \vskip0pt%
}

% Remove default navigation symbols
\setbeamertemplate{navigation symbols}{}

% --- METADATA ---
\title[Factor SV Models]{Efficient Bayesian Inference for Multivariate Factor Stochastic Volatility Models}
\subtitle{\citet{Kastner2017}}
\author{[Your Name]}
\institute[WU]{Vienna University of Business and Economics \\
        \medskip
        \scriptsize Bayesian Econometrics -- Summer Term 2026}
\date{[Presentation Date]}

%==== Beginning of the document ===============================================%

\IfFileExists{upquote.sty}{\usepackage{upquote}}{}
\begin{document}

%==== Title Page ==============================================================%

\begin{frame}
  \titlepage
\end{frame}

%==== Outline =================================================================%

\begin{frame}{Outline}
  \tableofcontents
\end{frame}

%==== 01 - Motivation =========================================================%

\section{Motivation}
\begin{frame}{From Univariate to Multivariate Stochastic Volatility}
  \begin{itemize}
    \setlength\itemsep{1em}

    \item \textbf{Recall Assignment Task 2:} We modelled daily CHF/USD returns
    with a $t_\nu$-distributed error and latent weights $\omega_i$, giving each
    observation its own effective variance $\sigma^2/\omega_i$ --- a
    \emph{univariate} stochastic volatility model.

    \item \textbf{Natural extension:} With $m$ assets we need not only
    time-varying variances but also time-varying \emph{covariances}.
    The full covariance matrix has $m(m+1)/2$ free elements --- growing
    \textbf{quadratically} in $m$.

    \item \textbf{The curse of dimensionality:} For $m = 26$ exchange rates
    this gives 351 covariance parameters \emph{per time point}. Direct
    modelling becomes infeasible.

    \item \textbf{Solution:} Impose a low-dimensional \textbf{factor structure}
    that reduces the parameter count while retaining flexibility for volatility
    clustering and co-movement.

  \end{itemize}
\end{frame}

%==== 02 - The Model ==========================================================%

\section{The Model}

%%% Slide 02A - Model Specification %%%
\begin{frame}{The Factor SV Model: Specification}

  For $t = 1, \ldots, T$ and $m$ observed returns $\mathbf{y}_t$ driven by
  $r \ll m$ latent factors $\mathbf{f}_t$:

  \begin{alertblock}{Observation \& Factor Equations}
    \[
      \mathbf{y}_t = \boldsymbol{\Lambda} \mathbf{f}_t + \boldsymbol{\varepsilon}_t,
      \qquad
      \mathbf{f}_t \mid \mathbf{h}_t \sim \mathcal{N}_r\!\left(\mathbf{0},\,
      V_t(\mathbf{h}_t^V)\right)
    \]
    with $\boldsymbol{\varepsilon}_t \mid \mathbf{h}_t \sim
    \mathcal{N}_m\!\left(\mathbf{0}, U_t(\mathbf{h}_t^U)\right)$
  \end{alertblock}

  \begin{itemize}
    \setlength\itemsep{0.7em}
    \item $\boldsymbol{\Lambda}$: unknown $m \times r$ \textbf{factor loadings} matrix
    \item $U_t = \mathrm{diag}\!\left(e^{h_{1t}}, \ldots, e^{h_{mt}}\right)$:
          \textbf{idiosyncratic} (series-specific) variances
    \item $V_t = \mathrm{diag}\!\left(e^{h_{m+1,t}}, \ldots, e^{h_{m+r,t}}\right)$:
          \textbf{factor} variances
    \item Each log-variance follows an independent \textbf{AR(1)} process:
          $h_{it} = \mu_i + \phi_i(h_{i,t-1} - \mu_i) + \sigma_i\eta_{it}$
  \end{itemize}

  \vspace{0.5em}

  The time-varying conditional covariance matrix is:
  \[
    \mathrm{cov}(\mathbf{y}_t \mid \mathbf{h}_t)
    = \boldsymbol{\Lambda} V_t \boldsymbol{\Lambda}^\prime + U_t
  \]

\end{frame}

%%% Slide 02B - Identification %%%
\begin{frame}{Identification}

  \begin{itemize}
    \setlength\itemsep{0.9em}

    \item \textbf{Scaling ambiguity:} Multiplying column $j$ of
    $\boldsymbol{\Lambda}$ by $c$ and dividing $f_{jt}$ by $c$ leaves
    $\mathbf{y}_t$ unchanged --- the model is not identified without
    additional constraints.

    \item \textbf{Classical fix:} Set $\lambda_{jj} = 1$ and treat the
    factor variance level as free.  This imposes a series ordering and creates
    ``leading'' factors.

    \item \textbf{Paper's approach:} Fix the level $\mu_{m+j} = 0$ for each
    factor log-variance (i.e.\ factor variance averages to 1), while leaving
    $\lambda_{jj}$ \emph{unrestricted}.

    \item \textbf{Sign switching:} Each column of $\boldsymbol{\Lambda}$ is
    identified only up to a sign flip --- resolved in post-processing.

  \end{itemize}

  \vspace{0.5em}

  \begin{block}{Why it matters}
    Unresolved non-identifiability causes MCMC chains to be multimodal
    and essentially useless.  The paper deliberately \emph{exploits} the
    non-identifiability as a computational tool (see ASIS, Section 4).
  \end{block}

\end{frame}

%==== 03 - Bayesian Inference =================================================%

\section{Bayesian Inference}

%%% Slide 03A - Gibbs Sampler %%%
\begin{frame}{Bayesian Inference: Gibbs Sampler}

  \textbf{Data augmentation} with latent volatilities $\mathbf{h}$ and factors
  $\mathbf{f}$ makes conditional posteriors tractable --- exactly as in the
  $t$-model from Task 2.

  \vspace{0.5em}

  \begin{alertblock}{Algorithm 1 -- Standard Gibbs Sampler}
    \begin{enumerate}
      \setlength\itemsep{0.5em}
      \item[\textbf{(a)}] \textbf{Univariate SV updates} ($m + r$ independent
            problems): given $\mathbf{f}$ and $\boldsymbol{\Lambda}$, each
            series/factor reduces to a univariate SV model via the
            log-squared observation trick.  Use the \texttt{stochvol} R
            package as a plug-in.
      \item[\textbf{(b)}] \textbf{Sample loadings} $\boldsymbol{\Lambda}$:
            conditional on $\mathbf{f}$ and $\mathbf{h}$, each row of
            $\boldsymbol{\Lambda}$ is a standard normal regression posterior.
      \item[\textbf{(c)}] \textbf{Sample factors} $\mathbf{f}_t$: conditional
            on $\boldsymbol{\Lambda}$ and $\mathbf{h}$, each $\mathbf{f}_t$
            is a multivariate normal posterior.
    \end{enumerate}
  \end{alertblock}

  \vspace{0.4em}

  \textbf{Link to lecture (Week 2--3):} Steps (b) and (c) are identical in
  structure to the conjugate Normal--Normal updates for $\mu$ and the
  $\omega_i$ in Task 2.  Step (a) extends the MH step for $\nu$.

\end{frame}

%%% Slide 03B - Why Standard Gibbs Fails %%%
\begin{frame}{Problem: Standard Gibbs Mixes Poorly}

  \begin{itemize}
    \setlength\itemsep{0.9em}

    \item The factor loadings $\boldsymbol{\Lambda}$ and the latent factors
    $\mathbf{f}_t$ are \textbf{strongly correlated} in their joint posterior.
    Updating them sequentially leaves the chain nearly stuck.

    \item In simulation studies with $m=10$, $r=2$, $T=1000$: the standard
    sampler yields \textbf{inefficiency factors (IFs) of 5000+} for some
    datasets --- meaning you need 5000 MCMC draws to equal one independent
    draw.

    \item Even after 5 million iterations the chain may not have converged.

  \end{itemize}

  \vspace{0.3em}

  \begin{block}{Inefficiency Factor (IF)}
    Ratio of the MCMC numerical variance to the variance from i.i.d.\ draws.
    An IF of $k$ means you need $k$ times as many MCMC draws to achieve the
    same accuracy as independent sampling.  Estimated via integrated
    autocorrelation time (\texttt{coda} package).
  \end{block}

\end{frame}

%==== 04 - ASIS ==============================================================%

\section{The ASIS Innovation}
\begin{frame}{Ancillarity-Sufficiency Interweaving Strategy (ASIS)}

  \textbf{Key idea:} Two \emph{mathematically equivalent} parameterisations
  are each efficient in different regions of the parameter space.
  Switching between them at every iteration combines the best of both worlds
  \citep{Yu2011}.

  \vspace{0.5em}

  \begin{columns}[T]
    \begin{column}{0.48\textwidth}
      \begin{alertblock}{Shallow Interweaving}
        \begin{itemize}
          \item Transform to alternative param.\ where $\lambda_{jj} = 1$
          \item Resample $\lambda_{jj}$ from a \textbf{Generalised Inverse
                Gaussian} (GIG) posterior (conjugate)
          \item Transform back
          \item Affects $\boldsymbol{\Lambda}$ and $\mathbf{f}$ only
          \item $\approx 2$--$8\times$ efficiency gain
        \end{itemize}
      \end{alertblock}
    \end{column}
    \begin{column}{0.48\textwidth}
      \begin{alertblock}{Deep Interweaving}
        \begin{itemize}
          \item Additionally shift log-volatility path: $h^*_{m+j,t} =
                h_{m+j,t} + 2\log|\lambda_{jj}|$
          \item Resample $\lambda_{jj}$ via
                \textbf{independence MH} in this extended space
          \item Affects $\boldsymbol{\Lambda}$, $\mathbf{f}$ \emph{and}
                $\mathbf{h}$
          \item Up to $\mathbf{400\times}$ efficiency gain
        \end{itemize}
      \end{alertblock}
    \end{column}
  \end{columns}

  \vspace{0.5em}

  \textbf{No tuning required:} The method is fully automatic --- no user
  adjustment of proposal variances (contrast with the $c_\nu$ tuning in
  Task 2).

\end{frame}

%==== 05 - Simulation =========================================================%

\section{Simulation Evidence}
\begin{frame}{Simulation Evidence}

  Setup: $m=10$ series, $r=2$ factors, $T=1000$ observations, 100 replications,
  5.1 million draws per run.

  \vspace{0.5em}

  \begin{table}[!h]
    \centering
    \fontsize{8pt}{10pt}\selectfont
    \begin{tabular}{lrrr}
      \toprule
      & \textbf{No interweaving} & \textbf{Shallow} & \textbf{Deep} \\
      \midrule
      \cellcolor{gray!10}{$\lambda_{11}$ (average IF)}
        & \cellcolor{gray!10}{$>1000$} & \cellcolor{gray!10}{$\approx 200$}
        & \cellcolor{gray!10}{$\approx 5$} \\
      Median IF across $\boldsymbol{\Lambda}$ & $\sim$400--5000 & $\sim$50--400 & $\sim$3--20 \\
      \cellcolor{gray!10}{Efficiency gain vs.\ no interweaving}
        & \cellcolor{gray!10}{---} & \cellcolor{gray!10}{$2$--$8\times$}
        & \cellcolor{gray!10}{up to $\mathbf{400\times}$} \\
      Factor $f_{1,T}$ (IF) & high & moderate & low \\
      \cellcolor{gray!10}{Factor log-var $h_{m+1,T}$ (IF)}
        & \cellcolor{gray!10}{high} & \cellcolor{gray!10}{high}
        & \cellcolor{gray!10}{substantially reduced} \\
      \bottomrule
      \multicolumn{4}{l}{\scriptsize Source: Tables 1--2 of \citet{Kastner2017}.
      IFs estimated via \texttt{coda}.}
    \end{tabular}
    \captionsetup{font=tiny}
    \caption{Average inefficiency factors for the factor loadings matrix.}
  \end{table}

  \vspace{0.3em}

  \begin{itemize}
    \item Deep interweaving trace plots show \textbf{immediate mixing}; the
    standard sampler barely moves even after millions of draws (Figure 1 of
    the paper).
    \item Log predictive Bayes factors (cf.\ Week 3) correctly identify $r=2$
    as the true number of factors.
  \end{itemize}

\end{frame}

%==== 06 - Application ========================================================%

\section{Exchange Rate Application}
\begin{frame}{Application: 26 EUR Exchange Rates (2005--2015)}

  \begin{itemize}
    \setlength\itemsep{0.7em}

    \item \textbf{Data:} $m=26$ daily EUR exchange rates, $T=2649$ log-returns
    (demeaned).  Includes outlier episodes: CHF spike Jan.\ 2015, RUB
    collapse Dec.\ 2014.

    \item \textbf{Model choice:} $r=4$ factors selected via log predictive
    Bayes factors (Week 3, slide 8).  500\,000 post-burn-in draws.

    \item \textbf{Interpretable factors:}
    \begin{itemize}
      \item \textbf{Factor 1:} USD/CNY-driven global dollar factor
      (USD, CNY, HKD load highly)
      \item \textbf{Factor 2:} Eastern European / EM currencies
      (ZAR, HUF, PLN, TRY; JPY loads \emph{negatively})
      \item \textbf{Factor 3:} Commodity currencies (AUD, NZD, CAD, ZAR)
      \item \textbf{Factor 4:} Tiger Cub economies (MYR, KRW, PHP, SGD)
    \end{itemize}

    \item \textbf{Time-varying correlations:} CNY and HKD nearly perfectly
    correlated with USD throughout; RUB--USD correlation fell from $\approx
    0.9$ to $\approx 0.4$ during 2008--2009.

    \item \textbf{Robustness:} Idiosyncratic SV components absorb outliers
    (CHF, RUB) without distorting the factor structure --- analogous to how
    $\omega_i$ absorbed outliers in Task 2.

  \end{itemize}

\end{frame}

%==== 07 - Takeaways ==========================================================%

\section{Takeaways \& Practical Relevance}
\begin{frame}{Takeaways \& Practical Relevance}

  \textbf{Theoretical connections to the lecture:}

  \begin{itemize}
    \setlength\itemsep{0.6em}

    \item \textbf{Data augmentation (Week 2--3):} The factor SV Gibbs sampler
    is a direct generalisation of the $\omega_i$-augmented sampler in Task 2.
    Both make conditionally intractable models tractable.

    \item \textbf{Model selection (Week 3):} Log predictive Bayes factors
    (Kass--Raftery scale) are used to select $r$ --- identical framework to
    the marginal likelihood comparisons in the lecture.

    \item \textbf{Shrinkage (Week 3):} The Gaussian prior on $\boldsymbol{\Lambda}$
    acts as a Ridge penalty, regularising small loadings and preventing
    spurious factors.

  \end{itemize}

  \vspace{0.5em}

  \textbf{Practical use:}

  \begin{itemize}
    \setlength\itemsep{0.4em}
    \item The \texttt{factorstochvol} R package implements the full algorithm.
    \item Deep interweaving adds negligible computational cost but yields
    orders-of-magnitude better mixing --- always prefer it.
    \item Even a correctly specified Gibbs sampler can be
    \textbf{practically useless} without good mixing.  ASIS is a general
    remedy applicable beyond this specific model.
  \end{itemize}

\end{frame}

%==== 08 - Bibliography =======================================================%

\section*{Bibliography}
\begin{frame}{Bibliography}

\bibliographystyle{plainnat}
\bibliography{references}

\end{frame}

%==== END =====================================================================%

\end{document}
